\documentclass[a4paper,12pt]{article}

\usepackage[utf8]{inputenc}
\usepackage[T2A]{fontenc}
\usepackage[russian]{babel}
\usepackage{amsmath}
\usepackage{amssymb}
\usepackage{mathtext}

\begin{document}

\textbf{Вариант 1}

В левой полости давление постоянно $p_0$, а в правой полости оно имеет постоянную
составляющую и гармонически изменяющуюся по времени составляющую $p_0 + p_1^+(\omega t)$.
Закон изменения давления на правом торце представим в виде
\begin{equation}
    p_1^+ = p_{1m}^+ f_p(\omega t), \quad f_p(\omega t) = \exp(i\omega t),
\end{equation}
где $p_{1m}^+$ — амплитуда пульсаций давления на торце канала; $\omega$ — частота пульсаций;
$f_p(\omega t)$ — закон изменения давления.

Введем в рассмотрение безразмерные переменные
\begin{equation}
\begin{aligned}
    &\psi = \delta_0/\ell \ll 1, \quad \lambda = w_m/\delta_0 \ll 1, \quad
     \Re = \delta_0^2 \omega / \nu, \quad \tau = \omega t, \quad
     \xi = x/\ell, \quad \zeta = z/\delta_0, \\
    &V_z = w_m \omega U_\zeta, \quad w = w_m W, \quad u = u_m U, \quad
     V_x = w_m \omega U_\xi / \psi, \\
    &p = p_0 + w_m \rho \nu \omega (\delta_0 \psi^2)^{-1} P, \quad
     p^+ = w_m \rho \nu \omega (\delta_0 \psi^2)^{-1} P^+.
\end{aligned}
\end{equation}
Здесь $p$ — давление; $\rho$, $\nu$ — плотность и кинематический коэффициент вязкости
жидкости; $V_x$, $V_z$ — проекции скорости движения жидкости на оси координат,
$w_m$ — амплитуда прогиба пластины; $W$ — безразмерный прогиб пластины;
$u_m$ — амплитуда продольного перемещения пластины; $U$ — безразмерное продольное
перемещение пластины; $\psi$, $\lambda$, $\Re$ — параметры, характеризующие задачу.

\end{document}
